\documentclass{ieeeaccess}
\usepackage{cite}
\usepackage{amsmath,amssymb,amsfonts,amsthm}
\usepackage{mathtools}
\usepackage{graphicx}
\graphicspath{{../results/}{./}}
\usepackage{textcomp}
\usepackage[utf8]{inputenc}
\usepackage[T1]{fontenc}
\usepackage{bm}
\usepackage{algorithm}
\usepackage{algorithmic}
\usepackage{booktabs}
\usepackage{xcolor}
\usepackage{color}
\usepackage{hyperref}
\hypersetup{hidelinks}
\usepackage{multirow}
\usepackage{makecell}
\usepackage{array}
\usepackage{url}
\usepackage{pifont}
\usepackage{siunitx}
\makeatletter
\AtBeginDocument{\DeclareMathVersion{bold}
\SetSymbolFont{operators}{bold}{T1}{times}{b}{n}
\SetSymbolFont{NewLetters}{bold}{T1}{times}{b}{it}
\SetMathAlphabet{\mathrm}{bold}{T1}{times}{b}{n}
\SetMathAlphabet{\mathit}{bold}{T1}{times}{b}{it}
\SetMathAlphabet{\mathbf}{bold}{T1}{times}{b}{n}
\SetMathAlphabet{\mathtt}{bold}{OT1}{pcr}{b}{n}
\SetSymbolFont{symbols}{bold}{OMS}{cmsy}{b}{n}
\renewcommand\boldmath{\@nomath\boldmath\mathversion{bold}}}
\makeatother

% ── Theorem-like environments ──────────────────────────────────────────
\newtheorem{proposition}{Proposition}
\newtheorem{remark}{Remark}
\newtheorem{lemma}{Lemma}
\renewenvironment{proof}{\noindent\textit{Proof.}}{\hfill$\square$\medskip}

% ═══════════════════════ MATH COMMANDS ═══════════════════
\newcommand{\R}{\mathbb{R}}
\newcommand{\SOg}{SO(3)}
\newcommand{\norm}[1]{\left\| #1 \right\|}
\DeclareMathOperator{\atan}{atan2}
\DeclareMathOperator{\Tr}{tr}
\DeclareMathOperator{\Logm}{Log}
\newcommand{\vect}[1]{\bm{#1}}
\newcommand{\mat}[1]{\mathbf{#1}}
\newcommand{\dt}{\Delta t}
\newcommand{\cmark}{\textcolor{black}{\ding{51}}}
\newcommand{\xmark}{\textcolor{black}{\ding{55}}}
% Filled from the N=5 trans/active batteries (scripts/batch_{trans,active}.sh).
\newcommand{\driftPlain}{7}    % plain-EKF hover drift over the transition [%] (7.0+-0.6)
\newcommand{\driftGated}{3}    % gated-EKF held error over the transition [%] (3.3+-0.8)
% Filled from the probe sweep (scripts/analyze_probe_sweep.py).
\newcommand{\probeOptShape}{a lateral tilt}
\newcommand{\probeOptErr}{2.5}
\newcommand{\probeFixErr}{8}

\begin{document}
\history{Date of publication xxxx 00, 0000, date of current version xxxx 00, 0000.}
\doi{10.1109/ACCESS.2026.0000000}

\title{Observability-Aware Virtual Force Sensing and Self-Calibration for Quadrotors}

\author{{Jos\'e Varela-Ald\'as}\authorrefmark{1},~\IEEEmembership{Senior Member, IEEE},
{Bryan~S.~Guevara}\authorrefmark{2},
{Ang\'elica Carri\'on}\authorrefmark{3},
and {Tiago Nascimento}\authorrefmark{2}}

\address[1]{Centro de Investigaci\'on MIST, Facultad de Ingenier\'ias,
Universidad Tecnol\'ogica Indoam\'erica, \mbox{Ambato 180103}, Ecuador
(e-mail: josevarela@uti.edu.ec)}

\address[2]{Laboratory of Systems and Robotics Engineering (LASER),
Center of Informatics (CI), Federal University of Para\'iba (UFPB),
Jo\~ao Pessoa, PB 58051-900, Brazil
(e-mail: bryansgue@laser.ci.ufpb.br; tiago@laser.ufpb.br)}

\address[3]{Faculty of Digital Engineering and Emerging Technologies,
International University of Ecuador UIDE, Quito 170411, Ecuador
(e-mail: anquitoca@uide.edu.ec)}

\tfootnote{The research was funded by Universidad Tecnol\'ogica Indoam\'erica.}

\markboth
{J. Varela-Ald\'as \headeretal:~Observability-Aware Virtual Force Sensing and Self-Calibration for Quadrotors}
{J. Varela-Ald\'as \headeretal:~Observability-Aware Virtual Force Sensing and Self-Calibration for Quadrotors}

\corresp{Corresponding author: Jos\'e Varela-Ald\'as (e-mail: josevarela@uti.edu.ec).}

\begin{abstract}
Many aerial tasks, such as flying in wind or carrying a payload, require the
external force on the airframe, yet quadrotors rarely carry a force sensor. We turn
the onboard inertial measurement unit into a \emph{virtual force sensor}: the
external specific force is read directly from the accelerometer residual. The
obstacle is calibration, since the reading needs the vehicle mass, which is
unobservable near hover, where it is indistinguishable from a vertical force; a
standard extended Kalman filter then corrupts the mass. We make the sensor
\emph{observability-aware and self-calibrating}. An online identifiability
measure, the Fisher information on the mass, gates the calibration and freezes
the mass when it cannot be seen; and when the flight lacks the excitation to
recalibrate, the controller injects a minimal probe that the same measure
shapes. In simulation the sensor reads the external force with a per-axis
correlation of \num{0.93} under continuous wind; the gate holds the mass
calibration to \SI{3}{\percent} through a windy hover where an ungated filter
drifts to \SI{7}{\percent}; the probe recovers the mass from a wrong prior that
hover alone cannot identify; and feeding the sensed force to a model-predictive
controller reduces the tracking error by \SI{34}{} to \SI{80}{\percent} under
wind, most where the wind dominates the dynamics. The sensor degrades gracefully
under injected accelerometer and thrust errors and runs in real time at
\SI{100}{\hertz}.
\end{abstract}

\begin{keywords}
Aerial robotics, force estimation, inertial sensing, observability,
self-calibration, virtual sensors.
\end{keywords}

\titlepgskip=-21pt

\maketitle

% ════════════════════ INTRODUCTION ════════════════════
\section{Introduction}\label{sec:intro}

\PARstart{T}{he} external force acting on a quadrotor, from wind, a payload, or physical
contact, governs how it must be flown, yet small aerial platforms rarely carry a
force or wrench sensor: the payload and power budget favor reusing the sensors
already on board. The accelerometer offers a route. It measures the
non-gravitational specific force directly, so once gravity and the commanded
thrust are removed, what remains is the external specific force. This turns the
inertial measurement unit (IMU) into a \emph{virtual force sensor}, with no added
hardware, and the reading is instantaneous: no double-differentiation of position,
no estimator lag.

The virtual sensor has one calibration parameter, the vehicle mass, which converts
the measured acceleration into the thrust contribution that must be subtracted.
The mass is rarely measured online, it is copied from a datasheet or changed by a
payload, and a few-percent error biases the reading. Worse, the mass is only
\emph{conditionally identifiable}: in near-hover flight it is indistinguishable
from a vertical external force, so an estimator that calibrates it unconditionally
corrupts the very reading it supports. A virtual force sensor for aerial use therefore needs
to know \emph{when} its calibration is reliable and to be able to recalibrate
itself when it is not. That is the problem we address.

Online parameter identification for aerial robots has been studied with
maximum-likelihood frameworks~\cite{burri2018system} and with recursive
regression on inertial signals that recover the mass and inertia under physical
consistency constraints~\cite{wuest2019online,wensing2017geometric}; these
methods need excited flight to identify the inertial parameters. External-force
estimation has a parallel line of work based on momentum
observers~\cite{tomic2017external,ruggiero2014impedance}, which estimate the
force on the airframe from the dynamics but are biased by errors in the same
parameters above. Moving-horizon estimation has been applied to quadrotor system
identification~\cite{liu2023nonlinear}, building on the constrained-MHE
theory~\cite{rao2003constrained}; the unmodeled aerodynamics that learned
residual models capture are themselves an active
area~\cite{bauersfeld2021neurobem,torrente2021data}. On the control side,
adaptive and incremental schemes reject disturbances without an explicit force
estimate~\cite{smeur2016adaptive}, and learning-based adaptive controllers fit a
disturbance representation online to reject wind in agile
flight~\cite{oconnell2022neuralfly}. These methods cancel the disturbance but do
not reason about \emph{when} the underlying parameters are identifiable.

Two themes run through this work. First, the mass is recoverable only under
sufficient excitation, a persistency-of-excitation condition that classical
optimal input design addresses by shaping the input to maximize the information
about the unknown~\cite{mehra1974optimal}. Freezing a parameter update when it is
not persistently excited is itself classical in adaptive estimation, in the form
of dead-zone and covariance-management schemes~\cite{smeur2016adaptive}, and
exciting the input to restore information is classical optimal input
design~\cite{mehra1974optimal}. Our contribution is not either element alone
but the signal that drives both: unlike a dead-zone keyed to a generic
excitation proxy, and unlike a momentum observer that estimates the force without
quantifying its identifiability, we gate on the closed-form Schur complement, which
(i) isolates exactly the mass--vertical-force coupling, (ii) is the marginal Fisher
information on the mass (its inverse is the Cram\'er--Rao bound), and (iii), in its
dimensionless form $\tilde\sigma$, doubles as the
controller's trigger to restore observability.
Second, the external force and the inertial parameters are coupled, so estimating
them jointly can bias both. The near-hover mass--vertical-force degeneracy we use
is a known structural fact in aerial estimation; existing methods, to our
knowledge, do not quantify it online to decide \emph{when} a parameter is
reliable or to \emph{act} so that it becomes observable. In near-hover flight
the mass and the vertical external force enter the dynamics through the same
channel and are not jointly observable, so an estimator that updates both
distributes the error between them. This is the property we exploit.

This paper builds an observability-aware, self-calibrating virtual force sensor.
The external force enters the accelerometer almost algebraically, so the sensor
reads it accurately once the mass is calibrated; the mass, by contrast, is only
conditionally identifiable and is confounded with a vertical force near hover. We
make the calibration observability-aware on two levels. An extended Kalman filter
(EKF) monitors an online identifiability measure and freezes the mass calibration
when the parameter is unobservable, protecting the reading during hover where a
naive filter drifts. When the mass must be recalibrated but the flight is too
weakly exciting to reveal it, the model-predictive controller injects a minimal probe that
restores identifiability and then stops, so the sensor never relies on a parameter
it cannot see. Fig.~\ref{fig:system} shows the resulting loop and the role of the
identifiability signal $\tilde\sigma$.

\begin{figure}[t]
\centering
\includegraphics[width=\columnwidth]{fig_system.png}
\caption{Architecture of the self-calibrating virtual force sensor: the
identifiability $\tilde\sigma$ gates the mass calibration and triggers the
active probe; the NMPC feeds $\hat{\vect{d}}$ forward.}
\label{fig:system}
\end{figure}

The contributions of this work are the following. The bare reading of the external
specific force from the accelerometer residual is a known momentum/IMU-residual
technique~\cite{tomic2017external,ruggiero2014impedance,smeur2016adaptive}; our
contribution is the observability-aware self-calibration layer that makes such a
virtual force sensor reliable on a vehicle whose calibration parameter is only
conditionally identifiable.
\begin{itemize}
\item We give the observability analysis of the sensor's calibration: the mass and
the vertical external force are not jointly observable in near-hover flight, and we
derive an online identifiability measure, the Schur complement of the force block
in the thrust--acceleration regressor Gramian, which is the marginal Fisher
information on the mass (its inverse is the Cram\'er--Rao bound) and quantifies when
the calibration is reliable.
\item We make the sensor \emph{self-calibrating}. The EKF gates the mass
calibration by this measure, freezing it when unobservable so the reading is
protected during hover; and when the mass must be recalibrated but the flight
lacks sufficient excitation, the controller injects a probe to restore
identifiability. The same
$\tilde\sigma$ \emph{shapes} the probe, since maximizing it selects the
most informative excitation within the probe family, so a single observability
scalar governs gating, triggering, and excitation design.
\end{itemize}
The resulting sensor is characterized and validated end to end in
software-in-the-loop: it matches the simulator ground truth to about
\SI{0.4}{\newton} RMS per axis (per-axis correlation \num{0.93} under continuous
wind), a moving-horizon estimator does not improve on it because the reading is
algebraically observable, and feeding the sensed force to a model-predictive
controller reduces the tracking error by \SI{34}{} to \SI{80}{\percent} under
wind, in real time at \SI{100}{\hertz}.

This is a methodology and feasibility study in software-in-the-loop. We exercise
the sensor against injected accelerometer bias, scale-factor, and thrust-mismatch
errors to probe the dominant real-sensor effects, but a full hardware demonstration,
on a real IMU and thrust map, or on a manipulator whose joint-torque readings give
the same observability structure, is the natural next step and is deliberately left
to a dedicated follow-up.

The rest of the paper is organized as follows. Section~\ref{sec:model}
presents the translational model and the measurement equations.
Section~\ref{sec:obs} analyzes observability and derives the hover ambiguity and
the identifiability measure. Section~\ref{sec:mhe} describes the
observability-aware EKF and its mass gate. Section~\ref{sec:nmpc} presents the
disturbance-rejecting controller and the active excitation.
Section~\ref{sec:setup} reports the experimental setup and
Section~\ref{sec:results} the results. Section~\ref{sec:conclusion} concludes.

% ════════════════════ SYSTEM MODEL ════════════════════
\section{System Model}\label{sec:model}

We model the quadrotor as a rigid body and estimate only its translational
motion. The attitude is measured by the onboard estimator and enters the model
as a known input, which removes the orientation states from the estimator and
keeps it well-conditioned.

\subsection{Frames and notation}
Let $\mathcal{F}_W=\{O_W,\vect{x}_W,\vect{y}_W,\vect{z}_W\}$ be the inertial
world frame and $\mathcal{F}_B=\{O_B,\vect{x}_B,\vect{y}_B,\vect{z}_B\}$ the body
frame attached to the center of mass (CoM). The position of $\mathcal{F}_B$ with
respect to $\mathcal{F}_W$ is $\vect{p}\in\R^3$ and its velocity is
$\vect{v}\in\R^3$. The orientation of $\mathcal{F}_B$ with respect to
$\mathcal{F}_W$ is the rotation matrix $\mat{R}(\vect{q})\in\SOg$, parameterized
by the unit quaternion
$\vect{q}=[q_0,\ \vect{q}_v^{\!\top}]^{\!\top}\in\mathbb{S}^3$, with
$\mathbb{S}^3\triangleq\{\vect{q}\in\mathbb{H}:\norm{\vect{q}}=1\}$ following the
convention of~\cite{sola2017quaternion}. The body rate $\vect{\omega}\in\R^3$ is
the angular velocity of $\mathcal{F}_B$ with respect to $\mathcal{F}_W$ expressed
in $\mathcal{F}_B$. We write $[\cdot]_\times$ for the skew-symmetric operator
that maps a vector to the matrix form of the cross product. The thrust acts
along $\vect{z}_B$, so its direction in $\mathcal{F}_W$ is
$\vect{a}\triangleq\mat{R}(\vect{q})\,\vect{e}_3$ with
$\vect{e}_3=[0,0,1]^{\!\top}$.

\subsection{Rigid-body dynamics}
The quadrotor is a six-degree-of-freedom rigid body actuated by a collective
thrust $T$ along the body $z$ axis and by an inner-loop body-rate controller. Its
state is the position $\vect{p}$, velocity $\vect{v}$, attitude $\vect{q}$, and
body rate $\vect{\omega}\in\R^3$. The dynamics are
\begin{align}
\dot{\vect{p}} &= \vect{v}, \label{eq:dp}\\
\dot{\vect{v}} &= -g\,\vect{e}_3 + \frac{T}{m}\,\mat{R}(\vect{q})\,\vect{e}_3
                  + \vect{d} - c\,\norm{\vect{v}}\,\vect{v},
                  \label{eq:dv}\\
\dot{\vect{q}} &= \tfrac{1}{2}\,\mat{\Omega}(\vect{\omega})\,\vect{q},
                  \label{eq:dq}\\
\dot{\vect{\omega}} &= \tfrac{1}{\tau_{rc}}\,(\vect{\omega}_{\text{cmd}}-\vect{\omega}),
                  \label{eq:dw}
\end{align}
where $g$ is the gravitational acceleration, $m\in\R_{>0}$ is the mass, and $c$
is a quadratic drag coefficient that lumps the dominant unmodeled
aerodynamics~\cite{sun2022comparative}. The attitude kinematics \eqref{eq:dq} use
the matrix
$\mat{\Omega}(\vect{\omega})=\left[\begin{smallmatrix}0 & -\vect{\omega}^{\!\top}\\[2pt] \vect{\omega} & -[\vect{\omega}]_\times\end{smallmatrix}\right]$,
following the convention of~\cite{sola2017quaternion,sola2018micro}. The body rate follows the
first-order response \eqref{eq:dw} to the commanded rate
$\vect{\omega}_{\text{cmd}}$, with time constant $\tau_{rc}$ for the inner-loop
rate controller, while the collective thrust $T$ is commanded directly. The term
$\vect{d}\in\R^3$, expressed in $\mathcal{F}_W$, is the \textit{external
specific force}, the net external force per unit mass acting at the center of
mass,
\begin{equation}
\vect{d}\triangleq \frac{1}{m}\,\vect{f}_e ,
\label{eq:dfdef}
\end{equation}
where $\vect{f}_e\in\R^3$ is the resultant of the forces that the nominal model
does not represent: aerodynamic effects beyond the drag term, wind, and any
added payload. Equivalently, $\vect{d}$ is the residual between the true
acceleration and the acceleration that gravity and thrust predict in
\eqref{eq:dv}; the experiments set the explicit drag coefficient $c=0$, so $\vect{d}$
also lumps the aerodynamic drag. We treat $\vect{d}$ as quasi-static over the
estimation window in near-hover flight, where the velocity-dependent drag is
negligible and an added mass and a vertical force become indistinguishable
(Section~\ref{sec:obs}); under fast flight $\vect{d}$ varies with velocity, which
the random-walk process model tracks within its bandwidth. Thus $\vect{d}$ is a
\emph{lumped specific-force residual}: the external force (wind, payload) plus the
lumped aerodynamic drag, the former dominating near hover and the latter growing
with speed. The force estimate is validated against the simulator's external-force
channel, and the residual drag is what the random-walk bandwidth limits at high
speed (Section~\ref{sec:results}). The quantities
estimated online are the mass $m$ and the external force $\vect{d}$, which enter
the translational dynamics \eqref{eq:dv}; the rate-loop constant $\tau_{rc}$ is a
nominal known constant of the inner loop, not estimated.

For estimation the attitude is treated as a known input: $\vect{q}$, and hence
the thrust axis $\vect{a}$, are taken from the accurate onboard estimate, so the
estimator integrates only the translational dynamics
\eqref{eq:dp}--\eqref{eq:dv}, with the thrust $T$ as a known input. This avoids the
quaternion states, with their re-normalization and hemisphere bookkeeping, and
keeps the estimator well-conditioned over the window.

\subsection{Measurement model}
The onboard sensors provide several signals at each step: the position and
velocity from odometry, the attitude $\vect{q}$ and the body rate
$\vect{\omega}$ from the inertial estimator and the gyroscope, and the specific
force from the accelerometer. The translational estimator uses two of them and
treats the attitude as an input. The first is the position--velocity
measurement
\begin{equation}
\vect{y} = [\,\vect{p}^{\!\top},\ \vect{v}^{\!\top}\,]^{\!\top}\in\R^6 ,
\label{eq:y}
\end{equation}
provided by the odometry, which is noisy in the simulator. The second is the
accelerometer, which measures the specific force in $\mathcal{F}_B$, the
non-gravitational acceleration sensed by the IMU~\cite{forster2015imu}; rotated into
$\mathcal{F}_W$ with the measured attitude it reads
\begin{equation}
\vect{s} = \mat{R}(\vect{q})\,\vect{a}_{\text{IMU}}
         = \frac{T}{m}\,\vect{a} + \vect{d} - c\,\norm{\vect{v}}\,\vect{v},
\label{eq:imu}
\end{equation}
the non-gravitational part of \eqref{eq:dv}. The attitude $\vect{q}$ enters as
the known input through $\vect{a}=\mat{R}(\vect{q})\,\vect{e}_3$ rather than as a
state to estimate, and the body rate $\vect{\omega}$ is not needed by the
translational estimator. Because the thrust $T$ and the attitude $\vect{q}$ are
known, \eqref{eq:imu} is an algebraic measurement of $\vect{d}$: the
accelerometer residual senses the external force directly, which turns the
estimator into a virtual force sensor.

% ════════════════════ OBSERVABILITY ════════════════════
\section{Observability and Excitation}\label{sec:obs}

The accelerometer measurement \eqref{eq:imu} relates the unknowns to known
signals, so we analyze when the mass and the external force can be recovered
from it. The drag term is a known function of the measured velocity, so we move
it to the left-hand side and define the compensated specific force
$\tilde{\vect{s}}\triangleq\vect{s}+c\,\norm{\vect{v}}\,\vect{v}$. With the
inverse mass $\beta\triangleq 1/m$, \eqref{eq:imu} becomes
\begin{equation}
\tilde{\vect{s}} = T\,\vect{a}\,\beta + \vect{d}
                 = \mat{\Phi}\,\vect{\varphi},
\qquad
\mat{\Phi} = [\,T\vect{a}\ \ \mat{I}_3\,],\quad
\vect{\varphi}=[\beta,\ \vect{d}^{\!\top}]^{\!\top},
\label{eq:regressor}
\end{equation}
which is linear in the unknown $\vect{\varphi}\in\R^4$ with regressor
$\mat{\Phi}\in\R^{3\times4}$. The thrust $T$ and the thrust axis $\vect{a}$ are
known, so identifiability of $\vect{\varphi}$ is the identifiability of
$(m,\vect{d})$.

The accelerometer is what makes \eqref{eq:regressor} algebraic. Without it, the
force $\vect{d}$ would enter only through the velocity derivative
$\dot{\vect{v}}$ in \eqref{eq:dv} and would have to be recovered by
differentiating the noisy odometry twice, which amplifies noise and adds a
delay. Reading $\vect{d}$ from the specific force \eqref{eq:imu} removes that
differentiation and yields an instantaneous algebraic measurement, which is the
basis of the virtual force sensor and decouples the analysis of $\vect{d}$ from
the integration of the trajectory.

The estimator augments the position and velocity with the mass and the force.
The position and velocity are measured by the odometry \eqref{eq:y} at every step,
so the recoverability of the parameters reduces to the \emph{identifiability} of
$(m,\vect{d})$ from the accelerometer model \eqref{eq:regressor} accumulated over
the window, a persistency-of-excitation condition on the regressor rather than a
full nonlinear observability test of the augmented state. This identifiability
determines when the mass and the force can be separated, and is what we analyze
and exploit.

A single sample gives three equations for four unknowns, so $\vect{\varphi}$ is
recovered only by accumulating samples. Over a window of $N$ samples the
information is the Gramian $\mat{G}=\sum_{k=1}^{N}\mat{\Phi}_k^{\!\top}\mat{\Phi}_k$,
which has the block form
\begin{equation}
\mat{G}=
\begin{bmatrix}
\sum_k T_k^2 & \sum_k T_k\vect{a}_k^{\!\top}\\[2pt]
\sum_k T_k\vect{a}_k & N\,\mat{I}_3
\end{bmatrix},
\label{eq:gramian}
\end{equation}
where $\norm{\vect{a}_k}=1$ was used. The force block $N\mat{I}_3$ is always
nonsingular, so given the mass $\vect{d}$ is recovered whenever the
window is nonempty; the joint degeneracy that remains is confined to the
mass--vertical-force pair at hover (Proposition~\ref{prop:hover}). The
mass enters only through $\beta$, which is identifiable if and only if the Schur
complement of the force block is positive,
\begin{equation}
\sigma \triangleq \sum_k T_k^2 - \tfrac{1}{N}\,\Big\|\textstyle\sum_k T_k\vect{a}_k\Big\|^2 > 0 .
\label{eq:schur}
\end{equation}
The scalar $\sigma$ is the dispersion of the vectors $T_k\vect{a}_k$ about their
mean; it is zero only when they are identical across the window, which is the
case at hover. Its dimension is thrust squared (newtons squared, not to be
confused with the window length $N$), so its
absolute scale grows with the thrust level of the maneuver. To obtain a
dimensionless, maneuver-comparable measure we normalize it by the thrust energy
in the window,
\begin{equation}
\tilde\sigma \triangleq \frac{\sigma}{\sum_k T_k^2}\in[0,1],
\label{eq:snorm}
\end{equation}
the relative dispersion of the thrust direction: $\tilde\sigma=0$ at hover and
$\tilde\sigma\!\to\!1$ when the thrust vector is maximally spread. Dividing out
the absolute thrust level removes the dominant source of the threshold's
maneuver dependence, so a single $\tilde\sigma_{\min}$ classifies the regime
across maneuvers far better than the raw $\sigma$, as Section~\ref{sec:results}
quantifies.

\begin{proposition}\label{prop:hover}
In stationary level hover, where $\vect{a}_k=\vect{e}_3$ and $T_k=\bar{T}$ are
constant over the window, the inverse mass $\beta$ is not identifiable
($\sigma=0$). Equivalently, only the combination $\bar{T}\beta+d_z$ is
observable, while $\beta$ and the vertical force $d_z$ are not separately
observable.
\end{proposition}
\begin{proof}
With $\vect{a}_k=\vect{e}_3$ and $T_k=\bar{T}$ we have $\sum_k T_k^2=N\bar{T}^2$
and $\sum_k T_k\vect{a}_k=N\bar{T}\vect{e}_3$, so
$\sigma = N\bar{T}^2-\tfrac{1}{N}\,N^2\bar{T}^2 = 0$, and $\beta$ is not
identifiable by \eqref{eq:schur}. The structure is explicit in
\eqref{eq:regressor}: the vertical row reads
$\tilde{s}_z=\bar{T}\beta+d_z$, a single equation in the two unknowns
$\beta$ and $d_z$, whereas the horizontal rows read
$\tilde{s}_{x,y}=d_{x,y}$ since $\vect{a}=\vect{e}_3$. Hence $d_x$ and $d_y$ are
observable while $\beta$ and $d_z$ collapse onto their sum.
\end{proof}

\begin{remark}
The ambiguity is confined to the vertical axis and to the pair $(m,d_z)$: a
mass error and a vertical force produce the same accelerometer residual at
hover. The joint ambiguity is confined to the mean thrust direction, which at
and near hover is vertical, so the horizontal components $d_x$ and $d_y$ remain
observable there, which is why the external force is recovered in the directions
where wind dominates.
\end{remark}

Away from hover the thrust axis tilts and its magnitude varies, so the vectors
$T_k\vect{a}_k$ disperse and $\sigma>0$; the window is then persistently exciting for
the mass. Since the Cram\'er--Rao lower bound on
the mass estimate scales with $\sigma^{-1}$ (derived in \eqref{eq:crb} below), a
more aggressive trajectory admits a better-conditioned estimate, while the force,
conditioned by the block
$N\mat{I}_3$ given the mass, is read with the same accuracy in every regime with
an accurate held mass. The mass is thus
observable under motion and degenerate only at hover, which
Section~\ref{sec:results} confirms by identifying it accurately across the range
of flight speeds.

The scalar $\sigma$ is therefore an online measure of how well the mass can be
identified, read directly from the known thrust and attitude over the window with
no extra sensing. It also has an estimation-theoretic meaning. With accelerometer
noise of variance $r_s$, the joint Fisher information of $(\beta,\vect{d})$ over
the window is $\mat{G}/r_s$ with $\mat{G}$ the Gramian \eqref{eq:gramian}. The
$(1,1)$ block of its inverse is the inverse Schur complement, $[\mat{G}^{-1}]_{11}=1/\sigma$,
so the \emph{marginal} Fisher information of the inverse mass $\beta$, after
profiling out the jointly estimated force $\vect{d}$, is $\sigma/r_s$, and the
Cram\'er--Rao bound on any unbiased estimate is
\begin{equation}
\mathrm{var}(\hat\beta)\ \geq\ r_s\,[\mat{G}^{-1}]_{11}\,=\,r_s\,\sigma^{-1},
\label{eq:crb}
\end{equation}
which diverges as $\sigma\to0$ at hover and tightens as the maneuver disperses
the thrust vector. The bound \eqref{eq:crb} is the batch least-squares CRB under
zero-mean accelerometer noise of isotropic covariance $r_s\mat{I}_3$ and a
quasi-static $\vect{d}$ over the window; it is the limit the recursive filter
approaches under excitation, used here as a relative identifiability measure rather
than an absolute uncertainty. Equation~\eqref{eq:crb} gives the \emph{raw} scalar $\sigma$ an
estimation-theoretic meaning: $\sigma/r_s$ is the marginal Fisher information on the
mass, so $\sigma$ is the inverse Cram\'er--Rao bound (up to the noise variance).
The bound is stated on the inverse mass $\beta$; since
$\mathrm{var}(\hat m)\approx m^4\,\mathrm{var}(\hat\beta)$ for small errors, the
bound on the mass itself scales with the same $\sigma^{-1}$.
For \emph{gating} we use its dimensionless normalization $\tilde\sigma$
\eqref{eq:snorm}, which divides out the thrust level so that a single threshold
$\tilde\sigma_{\min}$ holds across maneuvers; $\sigma$ carries the information
content and $\tilde\sigma$ carries the decision. We do not invert \eqref{eq:crb} to
set $\tilde\sigma_{\min}$ analytically; it is tuned with margin between the hover and
maneuver values of $\tilde\sigma$ (Section~\ref{sec:results}), and \eqref{eq:crb}
reports the uncertainty that any admissible threshold tolerates. The estimator uses
$\tilde\sigma$ to decide when to trust the mass, and the controller of
Section~\ref{sec:nmpc} uses it to decide when to excite.

Proposition~\ref{prop:hover} has a direct operational reading, summarized in
Table~\ref{tab:obs}. The external force is recoverable in every regime given an
accurate held mass, and is
recovered by any recursive estimator that uses the accelerometer residual; the
mass is recoverable only when the
flight excites it ($\sigma$ large), and is then identified jointly with the
force. This motivates the two elements of our method: an estimator that updates
the mass only when $\sigma$ is large (Section~\ref{sec:mhe}), and a controller
that creates the excitation when the mass must be identified but the flight
lacks the excitation to reveal it (Section~\ref{sec:nmpc}).

\begin{table}[t]
\centering
\caption{Observability regimes (Proposition~\ref{prop:hover}). The ``degraded''
$d_z$ entry is the joint case; with an accurate held mass it is recovered
(Section~\ref{sec:results}).}
\label{tab:obs}
\begin{tabular}{lcc}
\toprule
Quantity & Hover & Maneuver \\
\midrule
mass $m$        & unobservable ($\sigma\!\approx\!0$)         & observable ($\sigma$ large) \\
$d_x,\,d_y$     & observable                                  & observable \\
$d_z$           & degraded ($m$--$d_z$)                       & observable \\
\bottomrule
\end{tabular}
\end{table}

% ════════════════════ ESTIMATOR ════════════════════
\section{Observability-Aware Estimation}\label{sec:mhe}

We estimate the translational state and the parameters of
Section~\ref{sec:model} with an extended Kalman filter (EKF) whose state augments
the position and velocity with the mass and the external force,
\begin{equation}
\vect{x}=[\,\vect{p}^{\!\top},\ \vect{v}^{\!\top},\ m,\ \vect{d}^{\!\top}\,]^{\!\top}\in\R^{10}.
\label{eq:state}
\end{equation}
The known inputs at each step are the thrust $T$ and the thrust axis $\vect{a}$.
The position and velocity follow \eqref{eq:dp}--\eqref{eq:dv}, while the mass and
the external force have no dynamics of their own and evolve as random walks with
diagonal process covariance $\mat{Q}=\mathrm{diag}(q_p\mat{I}_3,q_v\mat{I}_3,q_m,q_d\mat{I}_3)$
(Table~\ref{tab:params}). Treating the
attitude as a known input keeps the filter free of quaternion states and their
log-map singularity; the body rate is not needed by the translational estimator.

\subsection{Prediction and update}
Each step propagates the dynamics, $\hat{\vect{x}}^-=\vect{f}(\hat{\vect{x}})$
and $\mat{P}^-=\mat{F}\mat{P}\mat{F}^{\!\top}+\mat{Q}$ with
$\mat{F}=\partial\vect{f}/\partial\vect{x}$, and applies two measurement updates.
The odometry \eqref{eq:y} is a linear update through
$\mat{H}=[\mat{I}_6\ \mat{0}]$ with covariance $\mat{R}$. The accelerometer
\eqref{eq:imu} is a nonlinear update with model
$\vect{g}(\vect{x})=(T/m)\vect{a}+\vect{d}-c\,\norm{\vect{v}}\,\vect{v}$,
Jacobian $\mat{H}_s=\partial\vect{g}/\partial\vect{x}$, and covariance
$\mat{R}_s$; this residual is the term that drives $\vect{d}$ and makes the
filter a virtual force sensor. Both updates use the standard EKF gain
$\mat{K}=\mat{P}^-\mat{H}^{\!\top}(\mat{H}\mat{P}^-\mat{H}^{\!\top}+\mat{R})^{-1}$,
and the mass and force are clipped to physical bounds after each step.

\subsection{The mass gate}\label{sec:gate}
A plain EKF updates every state at every step, including the mass when it is
unobservable. Under the hover ambiguity of Proposition~\ref{prop:hover} this
attributes a vertical force to a mass change and corrupts the estimate. The
observability-aware filter instead conditions the mass update on the
identifiability $\tilde\sigma$ of \eqref{eq:snorm}, computed online over a sliding
window from the known thrust and attitude. When the mass is not persistently
excited, the corresponding row of the Kalman gain is set to zero and its process
noise is suppressed,
\begin{equation}
\mat{K}_m\leftarrow\vect{0},\quad q_m\leftarrow 0
\qquad\text{while}\quad \tilde\sigma<\tilde\sigma_{\min}.
\label{eq:gate}
\end{equation}
With $q_m=0$ and the mass having no dynamics of its own ($F_{mm}=1$), the
prediction leaves the mass variance unchanged, $P_m^-=P_m$, so while the gate is
closed the mass neither drifts nor inflates: it is held exactly at its last
well-identified value with its uncertainty. The realized hold is therefore as good
as the gate's closure: in still hover the gate stays shut and the mass is held
exactly, whereas a windy hover intermittently re-excites $\tilde\sigma$ and
re-opens the gate, leaving the residual movement reported in
Section~\ref{sec:results}. Suppressing the process noise is what
makes this hold; were the variance left to grow, a transient spike in
$\tilde\sigma$ would re-open the update with a large gain and let a single sample
move the mass. To reject such spikes the gate re-opens the update only after
$\tilde\sigma$ has stayed above $\tilde\sigma_{\min}$ for a dwell time $t_d$
(Table~\ref{tab:params}), not on an isolated sample; this separates the sustained
excitation of a maneuver or probe from a momentary gust. The force update is
never gated, since, given the held mass, $\vect{d}$ is observable throughout
(Section~\ref{sec:obs}).

\subsection{Choice of estimator: EKF versus moving horizon}\label{sec:nomhe}
The accelerometer makes $\vect{d}$ an algebraic measurement \eqref{eq:regressor},
so the force needs no temporal smoothing to be recovered, and the mass, once
excited, is identified by the recursive update alone. A moving-horizon estimator
over the same state and measurements was implemented for comparison and did not
improve the estimates (Section~\ref{sec:results}): the window adds cost without a
benefit when the informative measurement is instantaneous and the recursive
filter already carries the full measurement history. We therefore retain the EKF
and place the contribution in the observability logic rather than in the
estimator class.

% ════════════════════ CONTROLLER ════════════════════
\section{Closed-Loop Use: Control and Active Calibration}\label{sec:nmpc}

The virtual sensor is used in closed loop by a nonlinear model-predictive
controller (NMPC) that serves two roles: it rejects the sensed disturbance by
feedforward (the application), and it injects the active probe that recalibrates
the mass when needed (the self-calibration of Section~\ref{sec:active}). We first
describe the controller, then the probe.

The estimated force closes the loop through this NMPC, which
tracks a reference trajectory while rejecting the
disturbance by feedforward. The controller keeps the full pose, so its state is
the position, velocity, attitude, and body rate,
$\vect{x}_c=[\vect{p}^{\!\top},\vect{v}^{\!\top},\vect{q}^{\!\top},\vect{\omega}^{\!\top}]^{\!\top}\in\R^{13}$,
and its input is the collective thrust and the commanded body rate,
$\vect{u}=[T,\,\vect{\omega}_{\text{cmd}}^{\!\top}]^{\!\top}\in\R^{4}$. The
prediction model is the rigid-body dynamics
\eqref{eq:dp}--\eqref{eq:dw} with the estimated mass and the estimated
disturbance injected,
\begin{equation}
\dot{\vect{v}} = -g\,\vect{e}_3 + \frac{T}{\hat{m}}\,\mat{R}(\vect{q})\,\vect{e}_3
                  + \hat{\vect{d}},
\label{eq:nmpc-dyn}
\end{equation}
so that $\hat{\vect{d}}$ enters as an additive specific-force feedforward that
cancels the external force; the rate loop keeps the form \eqref{eq:dw} with the
nominal time constant $\tau_{rc}$. The prediction model carries no explicit drag term; because the
estimator runs with $c=0$, the aerodynamic drag is lumped into $\hat{\vect{d}}$
and is therefore compensated through the same feedforward. The disturbance
estimate is smoothed by an exponential moving average before injection (the mass
estimate is smoothed likewise, with the constants of Table~\ref{tab:params}). In the
closed-loop experiments the online update is the force alone: the mass in the
prediction model is held at its identified value (here the known
\SI{1.05}{\kilogram}), so it is the feedforward $\hat{\vect{d}}$ that closes the
loop, while the mass calibration protects the force reading rather than
retuning the controller. Because the gate acts only on the mass row of the
filter, the force estimate of the gated and the plain filter coincide, and the
injected $\hat{\vect{d}}$ is that common reading.

\subsection{Cost and constraints}
Over a horizon of $N_c$ nodes the controller minimizes
\begin{align}
J=\ &\sum_{k=0}^{N_c-1}\Big(
   \norm{\vect{p}_k-\vect{p}_k^{\text{ref}}}^2_{\mat{Q}_p}
  +\norm{\vect{v}_k-\vect{v}_k^{\text{ref}}}^2_{\mat{Q}_v}\nonumber\\
  &\qquad +\norm{\vect{e}_q(\vect{q}_k)}^2_{\mat{Q}_a}
  +\norm{\vect{u}_k-\vect{u}_k^0}^2_{\mat{R}_u}\Big)
  +\norm{\vect{e}_N}^2_{\mat{Q}_N},
\label{eq:nmpc-cost}
\end{align}
where $\vect{p}^{\text{ref}}$ and $\vect{v}^{\text{ref}}$ are the reference
position and velocity,
$\vect{e}_q=\mathrm{Log}(\vect{q}^{-1}\!\otimes\vect{q}^{\text{ref}})\in\R^3$ is
the attitude error through the quaternion logarithm, with the standard
hemisphere correction, and the input is regularized around the hover
equilibrium $\vect{u}^0=[\hat{m}g,\,\vect{0}^{\!\top}]^{\!\top}$ so that a wrong
nominal mass does not bias the regularization. The velocity reference
$\vect{v}^{\text{ref}}$ is a feedforward term that removes the tracking lag at
high speed; $\mat{R}_u$ weights the input deviation. The terminal error
$\vect{e}_N$ stacks the position, velocity, and attitude errors at stage $N_c$,
weighted by $\mat{Q}_N=\mathrm{diag}(\mat{Q}_p,\mat{Q}_v,\mat{Q}_a)$. The inputs
are bounded by
$0\leq T\leq T_{\max}$ and $\norm{\vect{\omega}_{\text{cmd}}}_\infty\leq\omega_{\max}$,
with $T_{\max}=\SI{49}{\newton}$ ($\approx4.8mg$) and
$\omega_{\max}=\SI{20}{\radian\per\second}$.

The problem is solved with sequential quadratic programming under the real-time
iteration scheme, HPIPM, and a Gauss--Newton Hessian within
acados~\cite{verschueren2022acados}, sharing the framework with the estimator. It
uses $N_c=31$ nodes over a \SI{1.5}{\second} horizon at \SI{100}{\hertz}, with a
per-step solve time below \SI{1}{\milli\second}.

\subsection{Active excitation}\label{sec:active}
The gate of Section~\ref{sec:gate} protects a mass that has already been
identified, but it cannot create the information the flight does not provide:
in sustained hover the mass stays unobservable and the estimate is frozen at
the prior it started from. When an up-to-date mass is required and the
identifiability is low, the controller closes this gap by actively exciting the
unobservable direction. The same identifiability that gates the estimate now
\emph{shapes} the excitation: by \eqref{eq:crb} the Schur complement is the Fisher
information on the mass, so the probe that maximizes it (approximately, its
dimensionless form $\tilde\sigma$, since the probe perturbs the thrust energy
only to second order about $T\approx mg$) per unit deviation
minimizes the mass uncertainty. We do not solve the continuous
optimal-input-design problem~\cite{mehra1974optimal}; instead we maximize
$\tilde\sigma$ over a small two-parameter probe family by an empirical sweep, the
most informative probe within that family. Proposition~\ref{prop:hover}
identifies the two levers: a vertical bob of amplitude $\delta_z$ that modulates the
thrust magnitude $T$, and a lateral motion of amplitude $\delta_\ell$ that tilts the
thrust axis $\vect{a}$. Both disperse the regressor $T\vect{a}$, so the probe is
\begin{equation}
\vect{p}^{\text{ref}}\!\leftarrow\!\vect{p}^{\text{ref}}\!+\!
\big[\delta_\ell(\cos\omega_\ell t\!-\!1),\,\delta_\ell\sin\omega_\ell t,\,\delta_z\sin\omega_z t\big]^{\!\top},
\;\; \tilde\sigma<\tilde\sigma_{\min},
\label{eq:probe}
\end{equation}
with $\omega_\ell=2\pi f_\ell$ and $\omega_z=2\pi f_z$ (Table~\ref{tab:params}),
and we choose the $(\delta_z,\delta_\ell)$ mix that maximizes $\tilde\sigma$ subject to a
fixed amplitude budget $\delta_z^2+\delta_\ell^2\le\rho^2$, identified empirically in
Section~\ref{sec:results}. The
probe is triggered only while the mass is both needed and unobservable; once the
estimate has converged it is withdrawn and the gate again protects the mass. This
is the active counterpart of the gate.

\subsection{Remarks on closed-loop behavior}
The model error reaching the controller stays bounded. The uncompensated force is
the estimation residual $\vect{d}-\hat{\vect{d}}$, small because $\vect{d}$ is
observable throughout, and the mass enters only as the multiplicative thrust
scaling $\hat m/m$, whose error is bounded by \eqref{eq:crb} where the mass is
observable and frozen at its identification value where it is not. A formal
stability analysis of the coupled estimator--controller loop is left to future
work.

% ════════════════════ EXPERIMENTAL SETUP ════════════════════
\section{Experimental Setup}\label{sec:setup}

We evaluate the virtual sensor and its self-calibration in software-in-the-loop;
this section describes the simulation environment, the flights and disturbance, the
baselines, and the metrics.

\subsection{Simulation environment}
All experiments run in software-in-the-loop with the MuJoCo physics
engine~\cite{todorov2012mujoco}. The quadrotor has a mass of \SI{1.05}{\kilogram},
the sum of the body geometries in the model description, which we take as the
ground truth for the mass estimate. The odometry is published with injected
noise on position and velocity, so the estimator is exercised on noisy
measurements, and
the simulator exposes the true external force applied to the airframe on a
dedicated channel, which serves as ground truth for the force estimate. The
estimator, controller, and the scripts that reproduce every figure will be
released on publication, together with a data dictionary that maps each logged
column to its physical meaning and the source file for every figure and table. The
parameters are listed in Table~\ref{tab:params}.

\begin{table}[t]
\centering
\caption{Estimator and controller parameters.}
\label{tab:params}
\begin{tabular}{lll}
\toprule
Block & Parameter & Value \\
\midrule
Plant / sim       & mass, rate $1/\tau_{rc}$    & \SI{1.05}{\kilogram}, \SI{18}{\per\second} \\
Drag              & $c$ (estimator)            & $0$ (lumped into $\hat{\vect{d}}$) \\
Odometry noise    & injected on $p,v$          & MuJoCo, released config \\
Wind              & profile, amplitude         & step / sine, $\sim$\SI{2}{\newton} \\
EKF meas.\ var.   & $r_p,\,r_v,\,r_s$           & $0.02,\ 0.1,\ 0.5$ \\
EKF process       & $q_p,\,q_v,\,q_m,\,q_d$     & $10^{-6},10^{-3},10^{-5},10^{-2}$ \\
Mass gate         & window, dwell $t_d$, $\tilde\sigma_{\min}$ & \SI{1.0}{\second}, \SI{0.25}{\second}, $0.009$ \\
Smoothing         & $\alpha_m,\,\alpha_d$       & $0.995,\ 0.92$ \\
Active probe      & budget $\rho$, $f_\ell$, $f_z$, optimal & \SI{0.40}{\meter}, \SI{0.5}{\hertz}, \SI{1.0}{\hertz}, lateral \\
NMPC              & $N_c$, horizon             & $31$, \SI{1.5}{\second} \\
NMPC weights      & $\mat{Q}_p,\mat{Q}_a,\mat{Q}_v$ & $180,\ 20,\ 30$ \\
                  & $\mat{R}_u$                & $\mathrm{diag}(0.3,5,5,5)$ \\
NMPC limits       & $T_{\max},\,\omega_{\max}$  & \SI{49}{\newton}, \SI{20}{\radian\per\second} \\
\bottomrule
\end{tabular}
\end{table}

\subsection{Controller, trajectories, and disturbance}
The vehicle is flown by the NMPC of Section~\ref{sec:nmpc} at \SI{100}{\hertz},
the same rate as the estimator. Two references are used: a
stationary hover and a Lissajous figure-eight flown at peak speeds from
\SI{2.1}{} to \SI{10.2}{\meter\per\second}, which span the excitation levels of
the observability analysis. The external force is a \emph{deterministic} per-axis
step profile of about \SI{2}{\newton}, commanded by the controller on the
simulator's external-force channel; the simulator applies it and re-publishes the
\emph{applied} wrench as ground truth. Because the profile is defined in code, the
disturbance is identical across runs and its onset is controlled (it is armed only
once the vehicle has settled, after the verified reset), which makes every
experiment reproducible. The force estimate, by contrast, is read from the
accelerometer residual, independent of the commanded value, so the comparison
against ground truth is not circular. Sustained steps keep the disturbance within
the bandwidth of the momentum-based reading, examined in
Section~\ref{sec:results}. For the rejection study the same disturbance is applied
to the runs with and without the feedforward.

\subsection{Baseline}
The baseline is the \emph{plain EKF}, the same filter of Section~\ref{sec:mhe}
but without the mass gate, run on the identical measurement stream. This
ablation isolates the effect of the observability logic, since the two filters
differ only in the gate. As a check on the estimator class, we also ran a
moving-horizon estimator over the same state and measurements; it did not improve
the estimates (Section~\ref{sec:nomhe}) and is not carried as a separate
baseline.

\subsection{Metrics}
The virtual force sensor is assessed by the Pearson correlation between the
estimated force and the ground truth, per axis. Mass identification is assessed
by the error relative to the ground truth and its spread across repetitions.
Disturbance rejection is assessed by the position tracking error, reported as the
root-mean-square error with and without the feedforward under the same wind, and
its significance is tested with a one-sided Mann--Whitney $U$ test at each speed.
Each condition uses five repetitions with and five without the feedforward, so the
pooled estimation statistics are $N=10$ per speed; the hover feedforward case has
four (one run failed the reset check and was dropped), making it $N=4$ for the
hover rejection point and $N=9$ for the pooled hover estimation. We report
mean$\,\pm\,$std and the per-step solve time. The gate and active-excitation
figures (Figs.~\ref{fig:trans},~\ref{fig:active}) are each reported over five
repetitions with independent noise and wind realizations, shown as mean $\pm$ std;
the disturbance-free mass-identification figure (Fig.~\ref{fig:massconv}) is a
single representative run.

% ════════════════════ RESULTS ════════════════════
\section{Results}\label{sec:results}

We present the results in the order of the contributions: the external force, the
mass identification and its hover limit, the gate, the active excitation, the
disturbance rejection, robustness, and a summary comparison against both
baselines.

\subsection{The virtual force sensor}
The accelerometer makes the force algebraic \eqref{eq:regressor}, so the EKF
recovers it without any temporal smoothing. We report three complementary
measures so that no single one is misleading: the absolute error, the per-axis
Pearson correlation against the step ground truth (Fig.~\ref{fig:force}), and the
correlation under a continuously varying disturbance (Section~\ref{sec:robust}).
The absolute error is \SI{0.4}{\newton} root-mean-square per axis on a
$\sim$\SI{2}{\newton} disturbance, about \SI{20}{\percent} of the applied force;
the calibration curve below shows this residual is dominated by the bandwidth of
the random-walk force model, not by a scale error. The step correlation is
\num{0.91} to \num{0.98} across speeds (inflated by the step transitions, hence
reported alongside the absolute error); and under a sinusoidal wind the
correlation is \num{0.93} (Section~\ref{sec:robust}), the most demanding case.
The high-speed rows are also conservative in a specific sense: the ground-truth
channel contains only the applied wind, while the aerodynamic drag is
additionally lumped into $\hat{\vect{d}}$ ($c{=}0$ in the estimator), so the correlations at
speed mix the two; near hover, where the drag vanishes, the reading is validated
cleanly. A moving-horizon
estimator matches the EKF on these and does not improve them, confirming that the
force needs no special estimator beyond the accelerometer residual.
Fig.~\ref{fig:dhat} shows the estimate tracking the disturbance in the time
domain on all three axes.

The reading itself is the translational analogue of a classical
momentum/residual observer~\cite{tomic2017external,ruggiero2014impedance}, and we
do not claim to improve it: any such observer computed with the same mass would
return the same force. The comparison that matters is therefore not against
another force estimator but against the same estimator without the calibration
layer, which is the ablation of the following subsections;
Fig.~\ref{fig:mass} quantifies exactly the mass bias that an uncalibrated, or
unconditionally calibrated, observer inherits.

\begin{figure}[t]
\centering
\includegraphics[width=\columnwidth]{fig_dhat.png}
\caption{Virtual force sensor in the time domain (representative run):
$m\hat{\vect{d}}$ tracks the applied wind steps on all three axes.}
\label{fig:dhat}
\end{figure}

\begin{figure}[t]
\centering
\includegraphics[width=\columnwidth]{fig_force_corr.png}
\caption{Per-axis force correlation against excitation ($N{=}10$; $9$ in hover;
step ground truth).}
\label{fig:force}
\end{figure}

We characterize the sensor with a calibration curve, the standard test for a force
sensor: the estimated force ($m\hat{\vect{d}}$) against the applied force, per axis,
over a sinusoidal wind that sweeps the range (Fig.~\ref{fig:calib}). The fit is
linear and correctly scaled, with slopes \num{0.98}, \num{0.97}, \num{0.96} for the three axes,
sub-\SI{0.2}{\newton} intercepts, and $R^2$ of \num{0.90}, \num{0.89}, \num{0.87};
the unity slope confirms the reading is correctly scaled by the mass and the small
residual is the bandwidth limit of the random-walk force model. The sensor is
bandwidth-limited rather than range-limited: an amplitude sweep of step disturbances
shows the slope holding above \num{0.8} up to about \SI{2}{\newton} and then falling
(to \num{0.4} at \SI{4}{\newton}) as large, fast steps exceed the random-walk
bandwidth, so the usable full scale is set by how quickly the disturbance varies, a
limitation a higher-order force model would relax.

\begin{figure}[t]
\centering
\includegraphics[width=\columnwidth]{fig_calib.png}
\caption{Calibration curve: estimated force $m\hat{\vect{d}}$ versus applied
force, per axis, under a sinusoidal wind.}
\label{fig:calib}
\end{figure}

\subsection{Mass identification and its hover limit}
Given excitation and no confounding disturbance, the mass is identified
accurately. Fig.~\ref{fig:massconv} shows the EKF converging from a deliberately
wrong prior (\SI{0.60}{\kilogram}, \SI{43}{\percent} low) to within
\SI{0.6}{\percent} of the ground truth on a disturbance-free trajectory, with the
estimated force correctly near zero. This is the identification result: the mass
is a physical invariant and is recovered to the true value once the trajectory
excites the thrust direction.

\begin{figure}[t]
\centering
\includegraphics[width=\columnwidth]{fig_massconv.png}
\caption{Mass identification under excitation: from a \SI{0.60}{\kilogram} prior
to within \SI{0.6}{\percent} (disturbance-free trajectory).}
\label{fig:massconv}
\end{figure}

Hover is where identification is hard, and it exposes the limit of the ungated
filter. Fig.~\ref{fig:mass} shows the plain EKF mass against speed under wind
\emph{without} the gate: in hover the estimate is biased by \SI{4.8}{\percent} and
varies by $\pm\SI{35}{\gram}$ across the $N=9$ runs, because the vertical force is
misread as a mass change (Proposition~\ref{prop:hover}); under excitation the bias
falls below \SI{1}{\percent} and the spread to $\pm\SI{2}{\gram}$, in agreement
with the $\sigma^{-1}$ conditioning of Section~\ref{sec:obs}. This hover bias,
which appears even when the mass is seeded from the correct value, is exactly what
the gate prevents, as the next experiment shows.

\begin{figure}[t]
\centering
\includegraphics[width=\columnwidth]{fig_mass_id.png}
\caption{Mass observability against excitation, \emph{ungated} EKF ($N{=}10$ per
speed, $N{=}9$ in hover; mean $\pm$ std).}
\label{fig:mass}
\end{figure}

\subsection{The mass gate protects the estimate}
Fig.~\ref{fig:trans} flies an aggressive trajectory followed by a hover, both
under wind, over five runs with independent noise and wind. During the trajectory
both filters identify the mass. In the hover that follows, the plain EKF is blind
to the lost observability and attributes the vertical wind to a mass change; its
drift is more pronounced here than in the steady-hover battery of
Fig.~\ref{fig:mass} (\SI{4.8}{\percent}), because the maneuver leaves the filter
with an inflated mass covariance, so the wind moves the estimate by
\SI{\driftPlain}{\percent} over the transition ($7.0\pm0.6$, $N{=}5$). The
observability-aware filter detects $\tilde\sigma$ falling below
$\tilde\sigma_{\min}$ and gates the mass update, cutting the drift to
\SI{\driftGated}{\percent} ($3.3\pm0.8$); a residual movement remains because the
windy hover still excites $\tilde\sigma$ intermittently, but the gate more than
halves the error and is consistently better than the plain EKF across all runs.
The gate protects an estimate that hover can no longer fully support.

This advantage comes specifically from the Schur complement, not merely from gating
on excitation. We ran the same filter with the gate keyed instead to a generic
persistency-of-excitation proxy, the total regressor energy $\sum_k\norm{\mat{\Phi}_k}^2$,
a conventional dead-zone. Because hover is high in thrust energy though zero in
dispersion, this proxy never closes and the result coincides with the ungated plain
EKF (\SI{7.0}{\percent}): a generic excitation measure cannot tell hover from a
maneuver. Only $\tilde\sigma$, which measures the dispersion that isolates the
mass--vertical-force coupling, detects the unobservability and protects the
calibration.

\begin{figure}[t]
\centering
\includegraphics[width=\columnwidth]{fig_trans.png}
\caption{Mass over a windy maneuver-to-hover transition ($N{=}5$, mean $\pm$
std); hover shaded, $\tilde\sigma$ and threshold at the bottom.}
\label{fig:trans}
\end{figure}

\subsection{Active excitation recovers an unobservable mass}
When the mass is unknown and the vehicle must hover, no passive estimator can
recover it. Fig.~\ref{fig:active} starts from a deliberately wrong prior
(\SI{0.70}{\kilogram}) in hover. The passive filter stays frozen at the prior,
since hover provides no information (the passive estimate stays at
$-\SI{33.3}{\percent}$ across all $N{=}5$ runs); the active controller injects the
probe of \eqref{eq:probe}, $\tilde\sigma$ rises above the threshold, the gate
opens, and the mass converges from \SI{0.70}{} toward the ground truth, after which
the probe is withdrawn and the estimate is again protected. Using the
most informative probe (Section~\ref{sec:probeopt}), the mass error falls from
\SI{33}{\percent} to about \SI{3}{\percent} ($2.5\pm1.6$,
$N{=}5$); the recovery depends on the probe shape, which the next experiment
optimizes. By contrast, under a disturbance-free trajectory the same filter
identifies the mass to within \SI{0.6}{\percent} (Fig.~\ref{fig:massconv}); the
residual here is the price of identifying in hover from a wrong prior.

\begin{figure}[t]
\centering
\includegraphics[width=\columnwidth]{fig_active.png}
\caption{Mass recovery from a wrong hover prior (\SI{0.70}{\kilogram}; $N{=}5$,
mean $\pm$ std): the probe (shaded) lifts $\tilde\sigma$, recovers the mass,
then withdraws.}
\label{fig:active}
\end{figure}

\subsection{Shaping the probe by identifiability}\label{sec:probeopt}
The recovery in Fig.~\ref{fig:active} depends on the probe shape, and the
identifiability tells us which shape is best: since the Schur complement is the
Fisher information on the mass \eqref{eq:crb}, the probe that raises it most per unit
deviation tightens the bound. We sweep the vertical/lateral mix of \eqref{eq:probe}
at a fixed amplitude budget, from a pure lateral tilt to a pure vertical bob, and
measure the achieved $\tilde\sigma$ and the resulting mass recovery
(Fig.~\ref{fig:probeopt}, $N{=}5$ per shape). The lateral tilt raises $\tilde\sigma$
well above every other mix and is the only shape whose recovery is both accurate and
tightly distributed, consistent with \eqref{eq:crb}: it is a variance bound, so
the highest-$\tilde\sigma$ probe admits the lowest-variance recovery. The most informative
mix (\probeOptShape) recovers the mass to \SI{\probeOptErr}{\percent}, against
\SI{\probeFixErr}{\percent} for the original mixed shape at the same amplitude budget;
the difference is significant against every other shape (one-sided Mann--Whitney,
$p<0.01$). This shape is optimal only within the two-parameter family we sweep, not
the unconstrained optimal input.

\begin{figure}[t]
\centering
\includegraphics[width=\columnwidth]{fig_probe_opt.png}
\caption{Probe shaping: mass recovery error against the vertical/lateral mix at
a fixed amplitude budget; the $\tilde\sigma$-maximizing mix is circled.}
\label{fig:probeopt}
\end{figure}

\subsection{Disturbance rejection}
Feeding the estimated force to the controller reduces the tracking error under
wind, and the reduction scales with how much the wind dominates the dynamics
(Fig.~\ref{fig:rej}, Table~\ref{tab:rej}). Up to \SI{6}{\meter\per\second} the
wind dominates and the feedforward cuts the error by \SI{77}{} to
\SI{80}{\percent} (\SI{78}{\percent} in hover); as the speed grows the wind
becomes a small fraction of the aerodynamic forces and the reduction falls to
\SI{34}{\percent} at \SI{10}{\meter\per\second}. The improvement is statistically
significant at every speed: a one-sided Mann--Whitney $U$ test of the
feedforward against the baseline rejects the null at each operating point
($p<0.01$, $n=5$ per condition; $p=0.008$ for the $n=4$ hover case), with a large
pooled effect size (Cohen's $d=5.7$). The estimator
and controller each run at \SI{100}{\hertz} with a per-step time below
\SI{1}{\milli\second}. Fig.~\ref{fig:traj} shows the spatial effect on the
figure-eight: with the feedforward the vehicle stays on the reference, while
without it the wind pushes it off the path.

\begin{figure}[t]
\centering
\includegraphics[width=\columnwidth]{fig_traj.png}
\caption{Trajectory under wind: with the feedforward (blue) versus without
(red); reference dotted.}
\label{fig:traj}
\end{figure}

\begin{figure}[t]
\centering
\includegraphics[width=\columnwidth]{fig_rejection.png}
\caption{Disturbance rejection against speed ($N{=}5$ per mode; hover $N{=}4$;
mean $\pm$ std).}
\label{fig:rej}
\end{figure}

\begin{table}[t]
\centering
\caption{EKF estimation against trajectory speed ($N{=}10$; $9$ in hover): mass estimate and
per-axis force correlation with the ground truth. The force is read accurately
throughout; the (ungated) mass is biased and noisy in hover and accurate under
excitation.}
\label{tab:id}
\begin{tabular}{cccc}
\toprule
Peak speed [\si{\meter\per\second}] & $\hat m$ [\si{\kilogram}] & corr $d_x/d_y$ & corr $d_z$ \\
\midrule
hover & $1.101\pm0.035$ & $0.98/0.98$ & $0.97$ \\
2.1   & $1.083\pm0.027$ & $0.97/0.98$ & $0.98$ \\
3.8   & $1.049\pm0.028$ & $0.97/0.98$ & $0.95$ \\
6.1   & $1.051\pm0.010$ & $0.97/0.97$ & $0.91$ \\
8.2   & $1.056\pm0.003$ & $0.96/0.95$ & $0.94$ \\
10.2  & $1.053\pm0.002$ & $0.95/0.94$ & $0.93$ \\
\bottomrule
\end{tabular}
\end{table}

\begin{table}[t]
\centering
\caption{Disturbance rejection ($N=5$ per mode; hover $N=4$): tracking RMSE with and without the force
feedforward under the same wind, against peak speed.}
\label{tab:rej}
\begin{tabular}{cccc}
\toprule
Peak [\si{\meter\per\second}] & No ff [\si{\centi\meter}] & ff [\si{\centi\meter}] & Reduction \\
\midrule
hover & $38.7$ & $8.5$  & $78\%$ \\
2.1   & $38.0$ & $7.8$  & $80\%$ \\
3.8   & $40.0$ & $8.4$  & $79\%$ \\
6.1   & $36.0$ & $8.4$  & $77\%$ \\
8.2   & $32.7$ & $14.6$ & $55\%$ \\
10.2  & $32.8$ & $21.6$ & $34\%$ \\
\bottomrule
\end{tabular}
\end{table}

\subsection{Robustness}\label{sec:robust}
The virtual force sensor is not specific to the step disturbance used above.
Under a sinusoidal wind, a continuously varying profile, the estimate tracks the
ground truth with a per-axis correlation near \num{0.93} (Fig.~\ref{fig:sine}),
slightly below the step value because the moving disturbance has more bandwidth.
All results so far use the simulator's injected odometry noise, so the estimator
is already exercised on noisy measurements rather than a clean ground-truth
state. The gate threshold is stable across maneuvers once the identifiability is
normalized. The raw $\sigma$ separates the regimes but its absolute scale follows
the thrust level, so it required a maneuver-specific threshold (a span of more
than an order of magnitude across our experiments). The dimensionless
$\tilde\sigma$ of \eqref{eq:snorm} removes that dependence: settled hover sits at
$\tilde\sigma\approx0.005$ (windowed, peak $\approx0.007$) while an aggressive
trajectory reaches $\tilde\sigma\approx0.34$, nearly two orders of magnitude apart.
A single $\tilde\sigma_{\min}=0.009$ classifies hover, maneuver, and the active probe
alike: the probe lifts the instantaneous $\tilde\sigma$ above the threshold often
enough to re-open the gate (Fig.~\ref{fig:probeopt} reports the run median, which
sits just below it, with the lateral tilt the highest of the family), while hover
keeps the gate closed.
The remaining sensitivity is to transient gusts during a windy hover transition,
which can momentarily spike $\tilde\sigma$; the dwell time $t_d$, which requires the
threshold to be exceeded for a sustained interval, rejects these so the gate holds.
The threshold is thus dimensionless and shared across all regimes, with the dwell,
not the absolute signal scale, handling the residual gust sensitivity.

\begin{figure}[t]
\centering
\includegraphics[width=\columnwidth]{fig_dhat_sine.png}
\caption{Virtual force sensor under a sinusoidal wind (representative run);
per-axis correlation near \num{0.93}.}
\label{fig:sine}
\end{figure}

\subsection{Robustness to sensor errors}\label{sec:senserr}
A virtual force sensor lives or dies by the error sources a real IMU and thrust map
introduce, which the ideal simulator does not exercise. We inject them into the
estimator inputs and sweep their magnitude (Fig.~\ref{fig:robust}): an
accelerometer bias up to \SI{0.4}{\meter\per\square\second}, a scale-factor error
up to \SI{10}{\percent}, and a thrust-coefficient mismatch up to \SI{10}{\percent}.
The sensor degrades gracefully and proportionally, with no breakdown. The force
error rises from \SI{0.42}{\newton} to at most \SI{0.58}{\newton} (under the largest
bias), and the mass calibration from below \SI{1}{\percent} to at most
\SI{11.5}{\percent} (under a \SI{10}{\percent} thrust mismatch, which couples almost
one-to-one into the mass). The bias loads the force but barely the mass, while the
scale and thrust errors load the mass through the $T/m$ term, as expected; the
near-linear degradation indicates the calibration error is recoverable by a
one-time bench calibration of the bias and scale.

\begin{figure}[t]
\centering
\includegraphics[width=\columnwidth]{fig_robust.png}
\caption{Robustness to injected accelerometer bias, scale-factor error, and
thrust mismatch: force error (blue, left axis) and mass error (red, right
axis).}
\label{fig:robust}
\end{figure}

\subsection{Estimator comparison}
Table~\ref{tab:compare} consolidates the comparison on the identical data. The
two baselines fail in different ways, both traceable to the hover ambiguity of
Proposition~\ref{prop:hover}. The plain EKF reads the force well, including the
vertical axis, but lets the hover mass drift as the vertical wind is misread as a
mass change. The moving-horizon estimator is worse: its free mass both corrupts
the vertical force, which it misreads as mass, and drifts in hover (Fig.~\ref{fig:trans}),
and the window adds cost without recovering either. Our observability-aware EKF
keeps both: the EKF structure holds the vertical force, and the gate keeps the mass
closest to the truth in the windy hover (\SI{3.3}{\percent}, against
\SI{7.0}{\percent} for the plain EKF and \SI{5.3}{\percent} for the MHE). With
active excitation it is also the only one that recovers the mass from a wrong prior
in hover. All three run in real time at \SI{100}{\hertz}.

\begin{table}[t]
\centering
\caption{Estimator comparison on identical data. Force entries are pooled over
the battery; the excited-mass entries pool the well-excited runs
(8.2--10.2\,\si{\meter\per\second}, $N{=}10$ each speed); hover and recovery
entries are from Figs.~\ref{fig:trans} and~\ref{fig:active}.}
\label{tab:compare}
\begin{tabular}{lccc}
\toprule
 & MHE & plain EKF & obs-aware EKF \\
 &     & (baseline) & (ours) \\
\midrule
horizontal force corr      & 0.96 & 0.97 & 0.97 \\
vertical force corr        & 0.66 & 0.95 & 0.95 \\
mass error, excited        & $+0.4\%$ & $+0.4\%$ & $+0.4\%$ \\
mass error, windy hover    & $5.3\%$ & $7.0\%$ & $\mathbf{3.3\%}$ \\
recover mass, wrong prior  & no & no & yes \\
\bottomrule
\end{tabular}
\end{table}

% ════════════════════ CONCLUSION ════════════════════
\section{Conclusion}\label{sec:conclusion}

We built a self-calibrating virtual force sensor for quadrotors. The external
specific force is read directly from the accelerometer; its only calibration
parameter, the mass, is observable only under excitation and is confounded with a
vertical force in hover. We made the calibration observability-aware on two levels:
an extended Kalman filter that gates the mass calibration by the online
identifiability $\tilde\sigma$, freezing it when it cannot be seen, and a
model-predictive controller that injects a minimal probe to restore identifiability
when the mass must be recalibrated but the flight lacks the excitation. In
software-in-the-loop the sensor reads the external force to about \SI{0.4}{\newton}
RMS per axis (per-axis correlation \num{0.93} under continuous wind), the gate holds
the mass across a windy hover transition to \SI{3}{\percent} where a plain EKF drifts to
\SI{7}{\percent}, the active probe
recovers the mass from a wrong prior that hover alone cannot identify, and the
force feedforward reduces the tracking error by \SI{34}{} to \SI{80}{\percent}
under wind, scaling with how much the wind dominates the dynamics.

A matched comparison showed that a moving-horizon estimator does not improve on
the EKF for this problem, because the informative measurement is instantaneous;
the contribution therefore lies in the observability logic, not in the estimator
class. The dead-zone and the excitation are each classical; what ties them together
is the closed-form Schur complement, the Fisher information on the mass computed from
known signals (its inverse is the Cram\'er--Rao bound), whose dimensionless form
$\tilde\sigma$ drives the gate and the probe from one online scalar, including the
choice of the most informative excitation. The main
limitation is that the study is in simulation; hardware validation with a
motion-capture system and a calibrated disturbance is the immediate next step, as
is a continuous (rather than family-restricted) optimal-input design.
Because the filter reads any external specific force, the approach extends toward
contact and tether forces, where an observability-aware wrench estimate would
support aerial physical interaction without a dedicated force sensor.

% ════════════════════ REFERENCES ════════════════════
\begin{thebibliography}{99}

% --- External wrench / force estimation (momentum & residual observers) ---
\bibitem{tomic2017external}
T.~Tomi\'c, C.~Ott, S.~Haddadin,
``External Wrench Estimation, Collision Detection, and Reflex Reaction for Flying Robots,''
\emph{IEEE Trans.\ Robotics}, vol.~33, no.~6, pp.~1467--1482, 2017.

\bibitem{ruggiero2014impedance}
F.~Ruggiero \emph{et al.},
``Impedance Control of VToL UAVs with a Momentum-Based External Generalized Forces Estimator,''
in \emph{IEEE ICRA}, 2014, pp.~2093--2099.

% --- Moving-horizon estimation ---
\bibitem{rao2003constrained}
C.~V.~Rao, J.~B.~Rawlings, D.~Q.~Mayne,
``Constrained State Estimation for Nonlinear Discrete-Time Systems: Stability and Moving Horizon Approximations,''
\emph{IEEE Trans.\ Autom.\ Control}, vol.~48, no.~2, pp.~246--258, 2003.

\bibitem{liu2023nonlinear}
M.~Liu, Y.~Tan, I.~Mareels,
``Nonlinear Moving Horizon Estimation for Quadrotor System Identification,''
\emph{Control Eng.\ Practice}, vol.~135, art.~no.~105508, 2023.

% --- Online parameter / inertial identification ---
\bibitem{burri2018system}
M.~Burri \emph{et al.},
``A Framework for Maximum Likelihood Parameter Identification Applied on MAVs,''
\emph{J.\ Field Robotics}, vol.~35, no.~1, pp.~5--22, 2018.

\bibitem{wuest2019online}
V.~W\"uest, V.~Kumar, G.~Loianno,
``Online Estimation of Geometric and Inertia Parameters for Multirotor Aerial Vehicles,''
in \emph{IEEE ICRA}, 2019, pp.~1884--1890.

\bibitem{wensing2017geometric}
P.~M.~Wensing, S.~Kim, J.-J.~E.~Slotine,
``Linear Matrix Inequalities for Physically Consistent Inertial Parameter Identification,''
\emph{IEEE Robotics Autom.\ Lett.}, vol.~3, no.~1, pp.~60--67, 2018.

% --- Optimal input design / active sensing ---
\bibitem{mehra1974optimal}
R.~K.~Mehra,
``Optimal Input Signals for Parameter Estimation in Dynamic Systems: Survey and New Results,''
\emph{IEEE Trans.\ Autom.\ Control}, vol.~19, no.~6, pp.~753--768, 1974.

% --- Actuator dynamics / INDI ---
\bibitem{smeur2016adaptive}
E.~Smeur, Q.~Chu, G.~de Croon,
``Adaptive Incremental Nonlinear Dynamic Inversion for Attitude Control of Micro Air Vehicles,''
\emph{J.\ Guid.\ Control Dyn.}, vol.~39, no.~3, pp.~450--461, 2016.

\bibitem{bauersfeld2021neurobem}
L.~Bauersfeld, E.~Kaufmann, P.~Foehn, S.~Sun, D.~Scaramuzza,
``NeuroBEM: Hybrid Aerodynamic Quadrotor Model,''
in \emph{Robotics: Science and Systems}, 2021.

% --- Disturbance-aware / data-driven MPC for quadrotors ---
\bibitem{torrente2021data}
G.~Torrente, E.~Kaufmann, P.~Foehn, D.~Scaramuzza,
``Data-Driven MPC for Quadrotors,''
\emph{IEEE Robotics Autom.\ Lett.}, vol.~6, no.~2, pp.~3769--3776, 2021.

\bibitem{sun2022comparative}
S.~Sun, A.~Romero, P.~Foehn, E.~Kaufmann, D.~Scaramuzza,
``A Comparative Study of Nonlinear MPC and Differential-Flatness-Based Control for Quadrotor Agile Flight,''
\emph{IEEE Trans.\ Robotics}, vol.~38, no.~6, pp.~3357--3373, 2022.

\bibitem{oconnell2022neuralfly}
M.~O'Connell, G.~Shi, X.~Shi, K.~Azizzadenesheli, A.~Anandkumar, Y.~Yue, S.-J.~Chung,
``Neural-Fly Enables Rapid Learning for Agile Flight in Strong Winds,''
\emph{Science Robotics}, vol.~7, no.~66, art.~eabm6597, 2022.

% --- State-estimation foundations (notation / manifold) ---
\bibitem{sola2017quaternion}
J.~Sol\`a,
``Quaternion Kinematics for the Error-State Kalman Filter,''
arXiv:1711.02508, 2017.

\bibitem{sola2018micro}
J.~Sol\`a, J.~Deray, D.~Atchuthan,
``A Micro Lie Theory for State Estimation in Robotics,''
arXiv:1812.01537, 2018.

\bibitem{forster2015imu}
C.~Forster, L.~Carlone, F.~Dellaert, D.~Scaramuzza,
``IMU Preintegration on Manifold for Efficient Visual-Inertial Maximum-a-Posteriori Estimation,''
in \emph{Robotics: Science and Systems}, 2015.

% --- Tooling ---
\bibitem{verschueren2022acados}
R.~Verschueren \emph{et al.},
``acados: a modular open-source framework for fast embedded optimal control,''
\emph{Math.\ Program.\ Comput.}, vol.~14, no.~1, pp.~147--183, 2022.

\bibitem{todorov2012mujoco}
E.~Todorov, T.~Erez, Y.~Tassa,
``MuJoCo: A Physics Engine for Model-Based Control,''
in \emph{IEEE/RSJ IROS}, 2012, pp.~5026--5033.

\end{thebibliography}

\EOD

\vfill

\end{document}
